\begin{pastexam}[label={px:df3-106-ds-02}]{地方三等資訊處理 106 年 第二題}
（申論；20 分）

下表列出四種常見的資料結構，請填滿該表以顯示各資料結構在一般狀況下（average case），搜尋（search）、
插入（insertion）、刪除（deletion）資料之時間複雜度。陣列的各項資料已事先填入作為範例。（每小題 5 分，
共 20 分）
\begin{center}
\begin{tabular}{|l|c|c|c|}
\hline
 & \begin{tabular}{@{}c@{}}搜尋\\（search）\end{tabular} & \begin{tabular}{@{}c@{}}插入\\（insertion）\end{tabular} & \begin{tabular}{@{}c@{}}刪除\\（deletion）\end{tabular} \\
\hline
\multicolumn{1}{|c|}{陣列} & $O(n)$ & $O(n)$ & $O(n)$ \\
\hline
㈠佇列（queue） & & & \\
\hline
㈡雙向連結串列（doubly-linked list） & & & \\
\hline
㈢二元搜尋樹（binary search tree） & & & \\
\hline
㈣AVL樹（AVL tree） & & & \\
\hline
\end{tabular}
\end{center}
\end{pastexam}

\begin{pastsoln}
四列分屬資結書第~6、5、9、10~章，這裡歸在第~9~章。 $n$ 是資料筆數。 表中「插入」「刪除」照各資料
結構本身的操作來算：佇列是從尾端加入、從前端取出；串列是已經知道要在哪裡插入、要刪哪個節點。
\begin{center}
\begin{tabular}{lccc}
\toprule
 & 搜尋 & 插入 & 刪除 \\
\midrule
陣列（題目給的範例） & $O(n)$ & $O(n)$ & $O(n)$ \\
㈠ 佇列 & $O(n)$ & $O(1)$ & $O(1)$ \\
㈡ 雙向鏈結串列 & $O(n)$ & $O(1)$ & $O(1)$ \\
㈢ 二元搜尋樹 & $O(\log n)$ & $O(\log n)$ & $O(\log n)$ \\
㈣ AVL 樹 & $O(\log n)$ & $O(\log n)$ & $O(\log n)$ \\
\bottomrule
\end{tabular}
\end{center}

\textbf{㈠ 佇列。} 佇列沒有依鍵值排序，要找某筆資料只能從前端一個一個看，平均看一半，$O(n)$。
插入（enqueue）接在 rear、刪除（dequeue）從 front 拿，不管用環狀陣列還是鏈結串列，都只動幾個
指標或索引，$O(1)$（資結書第~6.2~節）。

\textbf{㈡ 雙向鏈結串列。} 搜尋要從頭沿著指標走，平均 $O(n)$。 插入：在已知節點的前或後（或串列
頭尾）接上新節點，只改四個指標，$O(1)$。 刪除：已知要刪的節點 $x$ 時，靠 $x.\mathit{prev}$ 直接
找到前一個節點，把前後兩個接起來，$O(1)$；這是雙向串列比單向串列好的地方，單向串列要先從頭找
前一個節點（資結書第~5.2~節）。 若「刪除某個鍵值」要先搜尋，就是 $O(n) + O(1) = O(n)$，作答時
可以補充這一點。

\textbf{㈢ 二元搜尋樹。} 三個操作都沿著一條從根往下的路徑，時間 $O(h)$，$h$ 是樹高。 鍵值以隨機
順序插入時，期望高度是 $O(\log n)$（資結書定理~9.2），所以一般狀況是 $O(\log n)$。 最差情況
（例如依序插入已排序的資料）樹退化成一條鏈，$h = n$，三個操作都是 $O(n)$。

\textbf{㈣ AVL 樹。} AVL 樹是高度平衡的二元搜尋樹，$n$ 個節點時高度一定是 $O(\log n)$（資結書
定理~10.1）。 搜尋走一條路徑；插入、刪除之後沿路往上調整平衡因子，必要時旋轉，每次旋轉 $O(1)$，
最多 $O(\log n)$ 次。 所以三個操作不只平均，最差情況也是 $O(\log n)$。
\end{pastsoln}
